\chapter{实验目的与任务说明}

\section{实验背景}

属性基加密（Attribute-Based Encryption, ABE）是一类将访问控制机制嵌入密码算法的公钥加密技术。与传统公钥加密“为某一个具体接收者加密”不同，ABE 允许加密方围绕属性集合或访问策略表达授权条件。例如，一个文件可以被加密给“部门为研发且安全级别不低于三级”的所有用户，而不是某一个固定身份的用户。这样的机制适合云存储、数据共享、组织权限管理和细粒度访问控制等场景。

本实验选择密文策略属性基加密（Ciphertext-Policy ABE, CP-ABE）作为实现对象。在 CP-ABE 中，密文携带访问策略，用户私钥携带属性集合。用户能否解密不由单一身份决定，而由其属性集合是否满足密文中的策略决定。本项目进一步选取线性秘密共享方案（Linear Secret Sharing Scheme, LSSS）作为访问策略的内部表示方法，将访问控制树转换为矩阵，并在解密端通过线性方程求解完成秘密重构。

因此，本实验并不是单纯实现一个“能加密、能解密”的命令行程序，而是围绕三个问题展开：第一，如何把直观的访问控制树转换为可计算的 LSSS 矩阵；第二，如何在 CP-ABE 的加密和解密过程中使用这些矩阵行；第三，如何通过实验说明满足策略和不满足策略时的行为差异。后续章节均围绕这三个问题展开，避免把算法、代码和实验结果割裂成互不相干的部分。

本报告的写作和实现均以 \verb|References_Article| 目录中的课程文件和论文为主要依据。课程期末文件将本题表述为“基于 LSSS 矩阵的属性基加密方案（使用访问控制树构建 LSSS 矩阵）”；实验三课件~[1]围绕 Waters 2011 方案~[3]与 FABEO 方案~[5]的实现与测试展开，给出了访问矩阵、属性映射、加密向量和解密重构常数等核心符号。因此，本项目把课程要求中的“访问控制树构建 LSSS 矩阵”作为实现主线，以实验三课件~[1]中的 Waters 2011 伪代码作为 CP-ABE 算法结构参照，并结合 Bethencourt、Sahai 和 Waters 提出的 CP-ABE 方案~[2]、Waters 的 LSSS-CP-ABE 构造~[3]以及 Liu、Cao 和 Wong 关于阈值访问树生成 LSSS 矩阵的研究~[4]补充理论依据。

\section{实验任务}

期末大作业要求在若干题目中任选一题完成完整实验报告和软件系统源码。本项目对应的题目是“基于 LSSS 矩阵的属性基加密方案（使用访问控制树构建 LSSS 矩阵）”。围绕这一要求，项目完成以下任务：

\begin{enumerate}
    \item 设计一种 JSON 格式的访问控制树表示方法，支持叶子属性、AND、OR 和一般阈值门。
    \item 实现从访问控制树到 LSSS 矩阵 $(M,\rho)$ 的转换算法。
    \item 实现 CP-ABE 的四个核心流程：Setup、KeyGen、Encrypt 和 Decrypt。
    \item 使用指数抽象模型验证 CP-ABE 的代数正确性。
    \item 在同一 \verb|Project| 工程中接入 BLS12-381 真实双线性配对库，验证同一访问结构在真实群元素上的可运行性。
    \item 使用 AES-CBC 与 HMAC-SHA256 完成明文载荷的对称加密和完整性校验。
    \item 提供命令行脚本和烟雾测试，使实验可以在 Windows PowerShell 环境中复现。
    \item 编写配套实验报告，解释理论基础、系统设计、实现细节和实验结果。
\end{enumerate}

上述任务中，最关键的是第二项、第三项和第五项。访问控制树到 LSSS 矩阵的转换决定了系统能否支持 AND、OR 和一般阈值门；Setup、KeyGen、Encrypt 和 Decrypt 四个流程则决定了矩阵结果能否真正参与属性基加密；真实配对库接入则用于检查底层群运算从指数抽象切换到真实 $G_1/G_2/G_T$ 元素后，访问树、矩阵构造和重构求解是否仍能保持稳定。对称加密和命令行脚本属于工程支撑部分，它们保证实验对象不是停留在公式层面，而是能够处理真实文件、生成可检查的 JSON 结果，并通过测试脚本反复验证。

\section{项目组织}

本仓库的主要目录如下：

\begin{table}[H]
\centering
\caption{项目目录组织}
\begin{tabular}{>{\raggedright\arraybackslash}p{0.30\textwidth} >{\raggedright\arraybackslash}p{0.60\textwidth}}
\toprule
目录或文件 & 作用 \\
\midrule
\verb|README.md| & 项目总说明，概述系统功能、目录结构、运行入口和实现边界。\\
\verb|Project/bin/lsss-abe.ps1| & 命令行入口，负责解析 setup、keygen、encrypt、decrypt 四类命令。\\
\path{Project/src/LsssAbe/LsssAbe.psm1} & 核心 PowerShell 模块，包含数学工具、LSSS 构造、加解密算法和 JSON I/O。\\
\verb|Project/examples| & 示例策略文件和示例明文。\\
\verb|Project/tests/smoke.ps1| & 烟雾测试脚本，验证九个成功、失败和边界输入场景。\\
\path{Project/src/LsssAbe/pairing_backend.py} & 真实 BLS12-381 pairing 后端，负责执行真实群元素、配对、哈希到群和序列化相关运算。\\
\path{Project/tests/smoke_pairing.ps1} & 真实 pairing 路线烟雾测试脚本，验证 BLS12-381 后端下的核心成功、失败和边界场景。\\
\verb|Report/LaTeX| & 本实验报告的 LaTeX 工程，包含正文、图表、截图和最终 PDF。\\
\verb|References_Article| & 课程要求、实验课件和相关论文。\\
\bottomrule
\end{tabular}
\end{table}

这样的目录组织体现了本实验的两条主线：一条是可运行的算法实现，另一条是对算法原理、实现细节和实验现象的解释。\verb|Project| 目录集中保存命令行入口、核心模块、示例文件、指数抽象路线和真实 pairing 路线；两条路线共用访问树表示、LSSS 矩阵构造、重构求解和测试组织方式，差异主要体现在底层群元素表示和配对运算实现。\verb|Report/LaTeX| 目录用于把访问控制树、LSSS 矩阵、线性重构、真实群元素实现和程序输出对应起来。参考资料单独保存在 \verb|References_Article| 中，则便于在论述算法来源和实验要求时保持材料来源清晰。

\section{实验目标}

本实验的直接目标不是完成可部署密码库，而是清晰展示“访问策略如何被转换为线性代数问题”，以及 CP-ABE 中“满足策略才能解密”的机制如何在代码中落地。因此，本项目将目标划分为理论目标、算法目标、工程目标和验证目标。

\begin{table}[H]
\centering
\caption{实验目标分解}
\begin{tabular}{p{0.22\textwidth} p{0.68\textwidth}}
\toprule
目标类别 & 具体内容 \\
\midrule
理论目标 & 理解 CP-ABE、访问控制树、阈值门、LSSS 矩阵和线性重构之间的关系。\\
算法目标 & 实现树到矩阵转换、高斯消元求重构系数、CP-ABE 四个核心算法。\\
工程目标 & 构建可运行命令行系统，使用 JSON 存储密钥和密文，便于检查和复现。\\
验证目标 & 通过正例和反例验证满足策略与不满足策略时的解密行为。\\
\bottomrule
\end{tabular}
\end{table}

这些目标之间存在递进关系。理论目标回答“为什么 LSSS 能表达访问策略”；算法目标回答“如何从树构造矩阵并求解重构系数”；工程目标回答“如何把算法组织成可运行的软件”；验证目标回答“实验结果是否与策略语义一致”。如果只完成前三类目标而缺少验证，系统只能说明“代码写出来了”；如果缺少理论解释，测试通过也难以说明其与 LSSS-ABE 题目的对应关系。因此本报告在结构上同时保留公式推导、代码节选和实验结果。

\section{运行环境}

本项目选择 Windows PowerShell 作为主要命令行环境。基础实现使用 .NET 标准库提供的 \verb|System.Numerics.BigInteger|、\verb|System.Security.Cryptography| 等组件完成大整数运算、随机数生成、哈希、AES 和 HMAC，从而在不依赖第三方密码库的情况下展示访问结构、线性重构和解密抵消关系。为了进一步检查底层群运算真实化后的系统行为，真实 pairing 路线额外使用仓库根目录下的 conda 环境 \verb|env|，由 Python 后端调用 \verb|py-ecc| 的 BLS12-381 群运算和 \verb|pairing| 接口，并使用 \verb|pycryptodome| 完成 AES-CBC/HMAC 载荷处理。

实验环境由四部分组成：第一部分是 PowerShell 脚本运行环境，用于执行 Setup、KeyGen、Encrypt、Decrypt 和烟雾测试；第二部分是 .NET 标准库，用于提供大整数和基础密码组件；第三部分是 Python pairing 后端，用于真实 BLS12-381 群元素、哈希到群、配对和群元素序列化；第四部分是 LaTeX 编译环境，用于生成最终实验报告。各部分的作用如表 \ref{tab:runtime-env} 所示。

\begin{table}[H]
\centering
\caption{实验运行环境组成}
\label{tab:runtime-env}
\begin{tabular}{>{\raggedright\arraybackslash}p{0.26\textwidth} >{\raggedright\arraybackslash}p{0.60\textwidth}}
\toprule
环境组成 & 作用 \\
\midrule
Windows PowerShell & 作为命令行入口和脚本执行环境，负责运行四个核心算法命令和测试脚本。\\
.NET 大整数与密码学库 & 提供 \verb|BigInteger|、随机数生成、SHA-256、AES、HMAC 等基础能力。\\
Python \verb|py-ecc| 与 \verb|pycryptodome| & 为真实 pairing 路线提供 BLS12-381 的 $G_1/G_2/G_T/\mathbb{Z}_p$、真实配对、哈希到群和对称载荷加密能力。\\
JSON 文件格式 & 保存公共参数、主密钥、用户私钥、密文、访问树和 LSSS 矩阵，便于实验检查。\\
XeLaTeX & 编译中文实验报告，支持公式、表格、代码清单和截图排版。\\
\bottomrule
\end{tabular}
\end{table}

真实 pairing 路线所使用的 Python 环境如图 \ref{fig:pairing-dependencies} 所示。截图中记录了 Python 版本、\verb|py-ecc| 版本和 \verb|pycryptodome| 版本，其中 \verb|py-ecc| 提供 BLS12-381 曲线和 pairing 运算，\verb|pycryptodome| 提供对称加密与消息认证组件。该环境证据与表 \ref{tab:runtime-env} 中的 Python pairing 后端相对应。

\begin{figure}[H]
\centering
\includegraphics[width=0.95\textwidth]{pairing_dependencies.png}
\caption{真实 pairing 路线依赖环境截图}
\label{fig:pairing-dependencies}
\end{figure}

从实验复现角度看，PowerShell 脚本使四个算法阶段具有清晰边界。Setup 生成公共参数和主密钥，KeyGen 根据属性集合生成私钥，Encrypt 把访问树和明文转换为密文 JSON，Decrypt 根据用户私钥和密文恢复明文或给出拒绝信息。每个阶段都会留下可检查的文件结果，后续章节中的代码说明、JSON 字段和终端截图都与这些文件相对应。

烟雾测试命令如下：

\begin{lstlisting}[language=bash,caption={烟雾测试命令}]
powershell -NoProfile -ExecutionPolicy Bypass -File .\Project\tests\smoke.ps1
\end{lstlisting}

真实 pairing 路线的验证命令如下：

\begin{lstlisting}[language=bash,caption={真实 pairing 版烟雾测试命令}]
powershell -NoProfile -ExecutionPolicy Bypass -File .\Project\tests\smoke_pairing.ps1
\end{lstlisting}

命令中使用 \verb|-ExecutionPolicy Bypass|，是因为部分 Windows 机器默认禁止直接执行 \verb|.ps1| 脚本。该参数只影响当前进程，不会永久改变系统策略，因而适合本实验的临时复现环境。

\section{报告结构}

围绕“访问控制树构建 LSSS 矩阵并用于属性基加密”这一主题，本文依次展开理论基础、方案设计、系统实现、实验结果和总结展望。这样的安排遵循从抽象概念到程序实现、再到实验验证的顺序：先交代访问结构和 LSSS 的数学含义，再说明这些概念如何进入代码、文件和测试结果。

第二章说明属性基加密、CP-ABE、访问控制树、LSSS、双线性对和对称加密之间的关系。该章的作用是建立术语和公式基础，尤其说明为什么访问策略可以被转化为矩阵重构问题。第三章给出系统方案设计，重点描述访问树到 LSSS 的构造方法，以及 Setup、KeyGen、Encrypt、Decrypt 四个核心算法之间的数据流。第四章进入系统实现层面，说明 PowerShell 模块结构、数据结构、矩阵转换函数、重构求解和 CLI 入口。

第五章是实验过程与结果分析，围绕九个指数版测试场景和八类真实 pairing 版补充场景说明系统在满足策略、不满足策略、一般阈值门、重复属性、特殊属性名、密文篡改和非法策略输入下的行为。第六章对实验完成情况、实现边界和真实安全证明条件进行总结，并说明指数抽象模型与 BLS12-381 真实配对实现之间的关系。附录则保留完整命令清单、示例策略 JSON、核心伪代码和参考文献，作为正文分析的补充材料。

整体上，第一章提出任务和环境，第二章提供理论依据，第三章给出方案设计，第四章呈现实现细节，第五章提供实验证据，第六章收束结论和边界。各章之间形成“题目要求—理论模型—算法设计—代码实现—实验验证—结论分析”的链路。
